\chapter{Pasywna enumeracja zasobów webowych i analiza aplikacji}\label{ch:domeny-web}

% ==============================================================================
% ROZDZIAŁ 4 - WSKAZÓWKI I TO-DO DLA AGENTA PISZĄCEGO:
% Rola: Osoba 2 (Aplikacje L7, subdomeny, serwery WWW, federacja edu-ID, archiwum i dorking).
% Cel: Identyfikacja punktów logowania SSO pod Evilginx2, taksonomia subdomen, procedury ISG.
% Narzędzia: crt.sh (wildcard SQL %), subfinder, amass, Wappalyzer, Wayback CDX, Google Dorking.
% Objętość docelowa: ~7-8 stron.
% ==============================================================================

\section{Architektura strefy domenowej i subdomenowej D-INFK}\label{sec:struktura-domenowa}

\subsection{Hierarchia przestrzeni nazw i domen powiązanych}\label{subsec:hierarchia-domen}
% TO-DO:
% 1. Opisać strukturę strefy inf.ethz.ch podpiętej pod nadrzędną ethz.ch.
% 2. Scharakteryzować 3 warstwy: centralna wydziału, instytuty badawcze (subdomeny lub domeny zewnętrzne
%    jak systems.ethz.ch, zisc.ethz.ch) oraz społeczność studencka (vis.ethz.ch).
% 3. Zwrócić uwagę na hybrydowy charakter strefy w kontekście wyznaczania granic rekonesansu.

\subsection{Pasywna enumeracja i taksonomia subdomen}\label{subsec:wyniki-enumeracji}
% TO-DO:
% 1. Przedstawić zapytanie do logów CT z symbolem wieloznacznym SQL LIKE (%): https://crt.sh/?q=%.inf.ethz.ch&output=json.
% 2. Pasywna agregacja subdomen z subfinder, amass (-passive) i crt.sh z deduplikacją unikalnych FQDN.
% 3. Syntetyczna taksonomia subdomen w 4 kategoriach (z podaniem przykładów np.):
%    - Bramy dostępowe i infrastruktura IT (np. isg, vpn, password, ssh).
%    - Usługi dydaktyczne i programistyczne (np. gitlab, moodle, svn, grader).
%    - Instytuty i laboratoria badawcze (np. spcl, las, netsec, systems).
%    - Klastry obliczeniowe i infrastruktura badawcza (z rozróżnieniem węzłów D-INFK od zasobów ogólnouczelnianych).

\section{Identyfikacja serwisów WWW i stosu technologicznego}\label{sec:systemy-webowe}

\subsection{Serwery WWW, systemy CMS i nagłówki HTTP}\label{subsec:technologie-cms}
% TO-DO:
% 1. Zidentyfikować system CMS na portalu głównym (Adobe Experience Manager - AEM, ścieżki /etc.clientlibs/ i /content/dam/).
% 2. Scharakteryzować serwery w laboratoriach (np. Nginx, Apache, generatory statyczne Hugo, Jekyll, GitLab Pages).
% 3. Tabela audytu nagłówków bezpieczeństwa HTTP (HSTS, CSP, X-Frame-Options, X-Content-Type-Options, Server) dla wybranych serwisów wydziałowych.
% 4. Zinterpretować wykryte luki w nagłówkach (np. brak CSP lub ujawnianie wersji oprogramowania) vs wzorcową konfigurację obronną (High Security Posture).

\subsection{Portale logowania i federacja tożsamości (edu-ID, AAI)}\label{subsec:portale-uwierzytelniania}
% TO-DO:
% 1. Opisać architekturę federacji tożsamości akademickiej: Switch edu-ID oraz Shibboleth AAI (SAML 2.0).
% 2. Zmapować dokładne punkty końcowe logowania (IdP): https://aai-logon.ethz.ch/... oraz https://eduid.ch/...
% 3. Wskazać punkty logowania GitLab (/users/sign_in) oraz portal samodzielnego zarządzania hasłami ISG (password.inf.ethz.ch).
% 4. Actionable Intelligence: powiązanie punktów logowania z atakiem Evilginx2 (przechwycenie ciasteczek sesyjnych
%    i ominięcie 2FA w Scenariuszu 1 w Rozdziale 6).

\section{Eksploracja archiwów internetowych}\label{sec:archiwa-internetowe}

\subsection{Archiwalne wersje stron w serwisie Wayback Machine}\label{subsec:wayback-machine}
% TO-DO:
% 1. Opisać metodykę zapytań do CDX Server API Wayback Machine:
%    https://web.archive.org/cdx/search/cdx?url=*.inf.ethz.ch/*&output=json&fl=original,timestamp,mimetype.
% 2. Przeanalizować etapy ewolucji: strony statyczne Solaris/Apache (1998-2006), Plone/Zope (2007-2015), konsolidacja na AEM.

\subsection{Wycofane podstrony i archiwalne artefakty techniczne}\label{subsec:wycofane-zasoby}
% TO-DO:
% 1. Wyszukanie w Wayback Machine historycznych artefaktów (np. dawne konfiguracje VPN/SSH, archiwalne ścieżki udziałów sieciowych NFS/SMB, nazewnictwo stacji studenckich).
% 2. Actionable Intelligence: wykorzystanie odnalezionych artefaktów do budowania wiarygodnych pretekstów technicznych w socjotechnice.

\section{Zaawansowane wyszukiwanie zasobów (Google Dorking)}\label{sec:google-dorking}

\subsection{Indeksowane katalogi i artefakty konfiguracyjne}\label{subsec:dorking-katalogi}
% TO-DO:
% 1. Tabela sprofilowanych operatorów Google Dorking dla domeny inf.ethz.ch (otwarte indeksy katalogów, dokumenty wewnętrzne, skrypty konfiguracyjne, komunikaty błędów).
% 2. Ocena perymetru: wykryte potencjalne ekspozycje vs potwierdzenie High Security Posture (brak wycieków kluczy API, tokenów czy plików środowiskowych .env).

\subsection{Publiczna dokumentacja operacyjna helpdesku ISG}\label{subsec:dorking-isg}
% TO-DO:
% 1. Przeanalizować bazę wiedzy ISG (isg.inf.ethz.ch) i otwarte wiki:
%    - Procedura onboardingu w CAB (pokój CAB E 41.1, odbiór karty ETH Card, przydział konta nethz).
%    - Podział stacji roboczych: Fully Managed Workstations vs Self-Administered (uprawnienia sudo).
%    - Oficjalne kanały wsparcia (helpdesk@inf.ethz.ch, telefon, system ticketowy).
% 2. Actionable Intelligence: wykorzystanie terminologii ISG, procedur i numerów pokoi jako legendy w scenariuszu vishingu w Rozdziale 6.
